\section{Separated rows and regeneration words}\label{sec:cuts}

We first construct independent successive pieces of a walk that stays
above its initial level in a prescribed direction.  Finite-range dependence requires a
gap between the rows used by consecutive pieces.  We create that gap by
marking some transitions whose probabilities are constant, so that the
walk can cross the gap without using its environmental rows.  The
construction will also prove the directional sign-union statement and a
finite mean for the spatial width of each piece.  Throughout this section
the direction may be any fixed nonzero real vector.  We adapt the
regeneration and record-count arguments of
\cite[Section~2]{OpenAI2026}, proving here the separated-row identities
needed for finite-range dependence.

\subsection{Marks and the rows a path uses}

We use the same splitting device as the independent-coin construction
of Comets and Zeitouni \cite[Equation~(2.1)]{CometsZeitouni2004}:
an environment-independent choice forces a step, and a residual choice
restores the prescribed transition row.  Two distinguished marks will
also prevent overlapping occurrences of the prescribed step sequence
introduced below.

Fix \(\lambda>0\) with \(2\lambda<\kappa\).  Enlarge the walk by giving
each step a mark in \(\{0,1,*\}\).  Conditional on the environment and
the marked past, a step from \(x\) has joint probabilities
\begin{equation}\label{eq:marked-kernel}
 \begin{aligned}
 P_{x,\omega}(X_1=x+e,\text{ mark }0)&=\lambda,\\
 P_{x,\omega}(X_1=x+e,\text{ mark }1)&=\lambda,\\
 P_{x,\omega}(X_1=x+e,\text{ mark }*)&=\omega(x,e)-2\lambda,
 \qquad e\in U.
 \end{aligned}
\end{equation}
These probabilities are nonnegative and sum to one.  Summing over the
marks recovers the original transition row, so every assertion about
the unmarked path is unchanged.  We keep the notation
\(P_{x,\omega}\) and \(P_x\) for the marked quenched and annealed laws.
We call these laws \emph{raw} when no condition of staying on one side of
a barrier is imposed.  The path filtration includes both positions and
the marks of steps already taken.  All constructions take place on the
full-measure set where the ellipticity inequalities hold at every site;
the marked kernel may be defined arbitrarily elsewhere.

A finite marked word \(\gamma\) consists of positions
\(x_0,\ldots,x_t\), with \(x_{j+1}-x_j\in U\), and one mark for each
of its \(t\) steps.  Its departure set and active departure set are
\[
 \Dep(\gamma)=\{x_0,\ldots,x_{t-1}\},\qquad
 \operatorname{Act}(\gamma)
 =\{x_j:0\le j<t,\ \text{step }j\text{ has mark }*\}.
\]
Let \(p_\omega(\gamma)\) be the product of the marked transition
probabilities in \eqref{eq:marked-kernel}, and put
\[
 W(\gamma)=E_Q[p_\omega(\gamma)].
\]
Thus \(W(\gamma)\) is the annealed probability of the specified prefix.
Stationarity makes it invariant under translation of all positions in
the word.  The terminal arrival uses no row.  More generally, a visited
row affects \(p_\omega(\gamma)\) only if it is an active departure.

\begin{lemma}[Separated path weights]\label{lem:separated-weights}
Let \(\gamma\) be a fixed finite marked word, and let
\(B\subseteq\ZZ^d\) be a deterministic set.  With the convention that
distance from an empty set is infinite, suppose that
\[
 \dist_\infty(\operatorname{Act}(\gamma),B)>R.
\]
If \(f\ge0\) is measurable with respect to the rows in \(B\), then
\begin{equation}\label{eq:separated-weights}
 E_Q[p_\omega(\gamma)f(\omega)]
 =W(\gamma)E_Q[f(\omega)].
\end{equation}
The conclusion also holds when the active set is empty, without a
separation requirement.

For a deterministic set \(F\subseteq\ZZ^d\), the quenched law of the
marked walk stopped on its first arrival in \(F\), including the arrival
site, depends only on rows outside \(F\).  In particular, the quenched
probability of any measurable infinite-path event that implies the walk
never enters \(F\) depends only on those rows.  Such probabilities may
therefore be used as \(f\) in \eqref{eq:separated-weights} whenever the
exterior rows are separated from \(\operatorname{Act}(\gamma)\).
\end{lemma}

\begin{proof}
Factors with marks \(0\) or \(1\) are constant.  All other factors in
\(p_\omega(\gamma)\) are functions of rows in
\(\operatorname{Act}(\gamma)\), with repeated departures from one row
retaining that same row.  The finite-range assumption
\eqref{eq:finite-range} gives independence from the rows in \(B\), also
when \(B\) is infinite.  This proves \eqref{eq:separated-weights}; if
the active set is empty, \(p_\omega(\gamma)\) itself is constant.

Before the first arrival in \(F\), every transition departs from its
complement.  Prefix probabilities for the stopped walk therefore use
only exterior rows, including when the prefix ends with that arrival.
These prefixes determine the stopped path law, so all its measurable
event probabilities have the same row dependence.  Equivalently, one
may kill the walk on arrival in \(F\).  An infinite-path event implying
avoidance of \(F\) is determined by this killed trajectory: it is false
on killed paths and retains its original meaning on surviving paths.
This proves the last assertion.
\end{proof}

We will always fix the relevant finite words before applying this
lemma.  The separated sets are then deterministic.  For example, two
prescribed words with departure sets at distance greater than \(R\)
have the same joint weight in independent environments as in a common
environment with conditionally independent walks.  The lemma also
applies to an avoidance event involving an entire infinite continuation.

\subsection{Cuts with independent continuations}

Fix a direction \(v\ne0\), rescale it so that
\(\norm{v}_\infty=1\), and write \(h(x)=v\cdot x\).  Choose a
coordinate step \(a_v\in U\) with \(h(a_v)=1\).  Fix an integer
\(M>dR+1\), the same integer for every direction considered, and set
\(\xi=\lambda^M\).  The prescribed \emph{script} in direction \(v\)
consists of \(M\) steps \(a_v\), the first with mark \(0\) and all
the others with mark \(1\).  Two occurrences of the script cannot
overlap in a step: the initial \(0\) of a second occurrence would have
to coincide with a \(1\) of the first.

We call time zero a record.  A positive time is an \emph{ordinary
record} if its height strictly exceeds all preceding heights.  A script
starting at a record has a \emph{candidate} at its end.  A candidate
time \(t\) is a \emph{cut} if
\[
 h(X_n)\ge h(X_t)\qquad(n\ge t).
\]
Each candidate is itself a strict record.  These definitions initially
refer to the whole path starting at time zero.  For a walk started
elsewhere we use heights relative to its starting site.  Write
\[
 D_v=\{h(X_n)\ge0\text{ for every }n\ge0\},\qquad
 p_v=P_0(D_v),\qquad
 \widehat P_v=P_0(\,\cdot\mid D_v)\quad(p_v>0).
\]

The script separates the rows of a record prefix from those of a
continuation above its endpoint; Figure~\ref{fig:cut-gap} displays this
distinction between visited and active sites.  Indeed, for all
\(x,y\in\ZZ^d\),
\begin{equation}\label{eq:height-distance}
 |h(x)-h(y)|\le d\norm{x-y}_\infty.
\end{equation}
If a record prefix ends at height \(r\), every departure in that
prefix has height strictly below \(r\).  The script has no active
departures, whereas a continuation staying above its endpoint uses
only rows of height at least \(r+M\).  The required separation follows
from \(M>dR\).

\begin{figure}[tb]
\centering
\begin{tikzpicture}[x=1cm,y=1cm,>=Stealth,font=\small]
 \fill[blue!7] (-3.6,-1.65) rectangle (3.7,0);
 \fill[green!8] (-3.6,2.15) rectangle (3.7,3.05);
 \draw[dashed,gray] (-3.6,0)--(3.7,0);
 \draw[dashed,gray] (-3.6,2.15)--(3.7,2.15);
 \draw[->,gray] (-4.05,-1.65)--(-4.05,3.15) node[above] {$h$};
 \node[left] at (-4.05,0) {$r$};
 \node[left] at (-4.05,2.15) {$r+M$};
 \draw[thick,blue!65!black]
   (-2.6,-1.25)--(-2.1,-.8)--(-1.5,-1.1)--(-1.1,-.5)
   --(-.6,-.75)--(0,0);
 \foreach \x/\y in {-2.6/-1.25,-1.5/-1.1,-1.1/-.5}
   \fill[blue!65!black] (\x,\y) circle (2pt);
 \foreach \x/\y in {-2.1/-.8,-.6/-.75}
   \draw[blue!65!black,fill=white] (\x,\y) circle (2pt);
 \foreach \y/\yy/\m in {0/.43/0,.43/.86/1,.86/1.29/1,1.29/1.72/1,1.72/2.15/1}
   \draw[->,thick,orange!85!black] (0,\y)--(0,\yy)
      node[midway,right=3pt] {\scriptsize $\m$};
 \foreach \y in {0,.43,.86,1.29,1.72,2.15}
   \draw[orange!85!black,fill=white] (0,\y) circle (2pt);
 \draw[thick,green!45!black,->]
   (0,2.15)--(.65,2.5)--(1.2,2.3)--(1.85,2.8)--(2.5,2.55);
 \node[align=left,anchor=west,blue!65!black] at (.9,-.75)
   {prefix active rows\\[-1pt]have $h<r$};
 \node[align=left,anchor=west,green!35!black] at (-3.4,2.6)
   {continuation rows\\[-1pt]have $h\ge r+M$};
 \node[align=left,anchor=west] at (.9,1.05)
   {script departures\\[-1pt]use constant weights};
 \draw[decorate,decoration={brace,amplitude=4pt}]
   (-.48,0)--(-.48,2.15) node[midway,left=7pt] {$M$};
\end{tikzpicture}
\caption{The gap at a candidate, shown schematically in height
coordinates.  Filled points indicate active prefix departures; orange
arrows are the constant-weight script steps.  The height gap
separates every active prefix row from every permitted continuation row
by more than \(R\) in \(\ell^\infty\)-distance, although the script
itself visits the intervening sites.}
\label{fig:cut-gap}
\end{figure}

Consequently, conditional on any finite marked prefix ending at a
record, the probability that the script occurs next and ends at a cut
is exactly
\begin{equation}\label{eq:script-probability}
 \xi p_v.
\end{equation}
More precisely, if \(\gamma_0\) is that prefix and \(B\) is any
event for a translated continuation from zero, then
\begin{equation}\label{eq:record-suffix-factorization}
 \begin{split}
 &P_0\{\text{prefix }\gamma_0,\ \text{then the script,}\\
 &\hspace{12mm}\text{then a translated continuation in }B\cap D_v\}\\
 &\hspace{12mm}=W(\gamma_0)\xi P_0(B\cap D_v).
 \end{split}
\end{equation}
Here and below a translated continuation records its positions
relative to its own starting site.  Lemma~\ref{lem:separated-weights}
proves this identity, using the forbidden halfspace below the script's
endpoint and stationarity.

We also record an elementary consequence of ellipticity that will be
used repeatedly.  A walk started in a slab
\(\{x:a\le h(x)\le b\}\), with \(a,b\) finite, exits that slab
almost surely under every uniformly elliptic quenched law.  Choose an
integer \(m>b-a\).  From any site of the slab, \(m\) consecutive
steps \(a_v\) leave it, and their probability is at least
\(\kappa^m\).  Conditional trials in successive blocks of \(m\)
steps bound the probability of remaining for \(jm\) steps by
\((1-\kappa^m)^j\).  In particular, on \(D_v\) the height supremum
is almost surely infinite: remaining forever in \([0,b]\) is null for
each positive integer \(b\).

To describe the pieces between cuts, call a finite marked word
\(\gamma\) from zero a \emph{regeneration word} if it has the
following properties:
\begin{enumerate}[label=(\roman*),leftmargin=*]
 \item all its heights are nonnegative;
 \item it ends at a candidate;
 \item every earlier candidate in the word is undercut later within
       the word, meaning that a subsequent height is strictly smaller
       than that candidate's height.
\end{enumerate}
Records and candidates in this definition are computed within the
word.  Its \emph{width} \(L(\gamma)\) is its terminal height.  The
last script starts at a record of height \(L(\gamma)-M\), so
\begin{equation}\label{eq:word-height-support}
 \begin{gathered}
 L(\gamma)\ge M,\qquad
 0\le h(x)<L(\gamma)\quad(x\in\Dep(\gamma)),\\
 h(x)<L(\gamma)-M\quad(x\in\operatorname{Act}(\gamma)).
 \end{gathered}
\end{equation}
We sometimes call a regeneration word a \emph{slab}, referring to this
height support.

\begin{proposition}[Regeneration law]\label{prop:cut-law}
Fix a normalized real direction \(v\) and suppose \(p_v>0\).
\begin{enumerate}[label=(\roman*),leftmargin=*]
 \item Under the raw law \(P_0\), on
       \(\{\sup_n h(X_n)=\infty\}\) there is a cut almost surely.
 \item Under \(\widehat P_v\), there are infinitely many cuts.  With
       time zero counted as an initial boundary, the relative marked
       words between successive boundaries are independent and
       identically distributed.  Their common law \(\nu_v\) assigns
       mass
       \begin{equation}\label{eq:regeneration-word-law}
        \nu_v(\{\gamma\})=W(\gamma)
       \end{equation}
       to each regeneration word \(\gamma\).  These finite words
       concatenate to the entire conditioned walk.
 \item Under the raw law, conditional on any realized prefix through
       the first cut, the translated suffix has law
       \(\widehat P_v\).  This assertion is restricted to the event
       that a cut exists.
\end{enumerate}
\end{proposition}

\begin{proof}
First test for a scripted cut from time zero.  Failure is detected in
finite time: either the next \(M\) steps are not the script, or they
are the script and the subsequent path first falls below its endpoint.
After a detected failure, wait for a strict record above the whole
maximum attained by the detection time, and test again.  The successive
test starts, when finite, are stopping times in the marked path
filtration.  Equation~\eqref{eq:script-probability}, applied to every
finite record prefix, gives conditional success probability
\(\xi p_v\) at each test.  Hence
\[
 P_0\{\text{the first }j\text{ tests occur and all fail}\}
 \le (1-\xi p_v)^j.
\]
On infinite height supremum, every detected failure is followed by
another finite test start.  Letting \(j\) tend to infinity proves (i).
The slab-exit observation implies that the first cut is finite almost
surely under \(\widehat P_v\).  The tests need not locate the earliest
cut; their only role is to establish that some cut exists.

We now identify the law at the actual first cut by finite prefixes.
Suppose \(\gamma\) is a
regeneration word ending at \(z\).  A path beginning with \(\gamma\)
whose suffix stays at or above \(h(z)\) has its first cut at that
endpoint: all earlier candidates have already been undercut.  It also
satisfies the original \(D_v\).  Conversely, a path in \(D_v\)
through its first cut gives such a word.  Indeed, every earlier
candidate has smaller height than the terminal strict record.  If it
were not undercut before the terminal record, it could not be undercut
afterwards either, since the suffix stays above that record; it would
already be a cut.

For any measurable suffix event \(B\), the row separation in
\eqref{eq:word-height-support} and
Lemma~\ref{lem:separated-weights} therefore give
\[
 \begin{split}
 &P_0\{\text{first-cut prefix }\gamma,\
          \text{ translated suffix in }B\}\\
 &\hspace{25mm}=W(\gamma)P_0(B\cap D_v).
 \end{split}
\]
Dividing by \(p_v\) proves that, under \(\widehat P_v\), the
first word has mass \(W(\gamma)\) and its translated suffix has law
\(\widehat P_v\), independently of that word.  There are only
countably many finite marked words, and the first cut is finite almost
surely, so these masses sum to one.

The cut is a strict record, hence records of its translated suffix are
also records of the full path.  A script cannot straddle this boundary:
such a script would overlap the script that has just ended.  Thus the
cuts of the suffix are precisely the subsequent cuts of the full path.
Iterating the preceding factorization proves the iid assertion in
(ii).  Every word is finite and contains at least \(M\) steps, so the
successive cut times tend to infinity and the words cover the whole
path.

For (iii), the prefix through the raw first cut may have negative
heights, but it still ends with the script from a record, and
every preceding candidate is undercut within the prefix.  Its active
departures lie strictly below the start of the terminal script.  The
same separated-weight calculation gives raw prefix probability
\(W(\gamma)p_v\) and, with suffix restriction \(B\), probability
\(W(\gamma)P_0(B\cap D_v)\).  Their ratio is
\(\widehat P_v(B)\).  All of these arguments enumerate finite
prefixes; none uses a Markov property at a cut selected using the
future.
\end{proof}

\subsection{Escape signs and finite mean width}

The regeneration law gives two consequences needed in the remainder of
the proof.  First, a positive escape probability in one direction rules
out oscillation between arbitrarily high and low levels.  Second, the
widths of the regeneration words have finite mean.  Counting records,
rather than integer height levels, will give the latter conclusion for
real directions as well.

\begin{proposition}[Directional sign union]\label{prop:sign-union}
For every fixed \(v\ne0\),
\begin{equation}\label{eq:sign-union}
 P_0(A_v)>0\quad\Longrightarrow\quad
 p_v>0\ \text{ and }\ P_0(A_v\cup A_{-v})=1.
\end{equation}
More precisely, when \(p_v>0\), infinite height supremum implies
\(A_v\) almost surely, and finite height supremum implies \(A_{-v}\)
almost surely.
\end{proposition}

\begin{proof}
Assume first that \(p_v>0\).  On infinite supremum there is a first
cut, and its suffix has the iid word law of
Proposition~\ref{prop:cut-law}.  The successive slab bases increase in
height by at least \(M\); within each slab the path stays above its
base.  Since each slab is finite, \(h(X_n)\to+\infty\).

For finite supremum, fix integers \(m,b\ge0\).  Consider the event
that all heights are at most \(m\), but heights at least \(-b\) are
visited infinitely often.  From each such visit, a script of
\(m+b+1\) consecutive steps \(a_v\), with no mark restriction,
crosses \(m\) with conditional probability at least
\(\kappa^{m+b+1}\).  Take successive visit trials at least
\(m+b+1\) times apart.  The probability that the first \(j\) trials
occur and all fail to cross \(m\) is at most
\((1-\kappa^{m+b+1})^j\).  The event in question is therefore null.
Taking the countable union over \(m,b\) shows that finite supremum
forces \(h(X_n)\to-\infty\).

It remains to obtain \(p_v>0\) from a positive probability of
\(A_v\).  On \(A_v\), the height sequence tends to infinity, so its
global minimum is attained at some finite time, even if its possible
values are not discrete.  Thus for some deterministic \(n\ge0\)
and \(x\in\ZZ^d\), the event
\[
 \{X_n=x,\ h(X_{n+j})\ge h(x)\text{ for all }j\ge0\}
\]
has positive annealed probability.  If \(q_x(\omega)\) denotes the
quenched probability from \(x\) of staying above \(h(x)\), the
quenched Markov property at time \(n\) expresses this probability as
\(E_Q[P_{0,\omega}(X_n=x)q_x(\omega)]\).  It is bounded above by
\(E_Qq_x=p_v\), by stationarity.  Therefore \(p_v>0\), and the
two alternatives just proved give \eqref{eq:sign-union}.
\end{proof}

\begin{proposition}[Finite mean spatial width]\label{prop:mean-width}
For every normalized real direction \(v\) with \(p_v>0\), the
regeneration law of Proposition~\ref{prop:cut-law} satisfies
\begin{equation}\label{eq:finite-mean-width}
 E_{\nu_v}L<\infty.
\end{equation}
\end{proposition}

\begin{proof}
Let \(S(\gamma)\) be the number of ordinary records in a regeneration
word \(\gamma\), excluding its initial boundary and including its
terminal record.  Each increase of the running maximum is at most one,
because every step has height increment at most one.  Therefore
\(L(\gamma)\le S(\gamma)\).  The record indices of successive cuts,
with index zero at the initial boundary, are sums of iid positive
integer increments with law \(S\): the terminal height of each word
is the maximum so far, so record counts concatenate across boundaries.

For \(i\ge0\), let \(\rho_i\) be the time of ordinary record
index \(i\) in a raw walk, with \(\rho_0=0\) and
\(\rho_i=\infty\) if that record is never reached.  Put
\[
 D'_{v,i}=\{\rho_i<\infty,\ h(X_j)\ge0
                       \text{ for }0\le j\le\rho_i\}.
\]
A script from record \(i\) creates exactly \(M\) consecutive
strict records.  Hence a cut of record index \(i+M\), on a path in
\(D_v\), occurs exactly when the prefix realizes \(D'_{v,i}\),
the script follows, and its suffix stays above its endpoint.  Summing
\eqref{eq:record-suffix-factorization} over the possible record
prefixes and dividing by \(p_v\) gives
\begin{equation}\label{eq:record-renewal-lower}
 \widehat P_v\{i+M\text{ is a cut index}\}
   =\xi P_0(D'_{v,i})\ge\xi p_v.
\end{equation}
The inequality holds because on \(D_v\) the supremum is infinite
almost surely, so every ordinary record is reached.

Let \(N_{\mathrm{rec}}(n)\) count the cut indices in \(\{1,\ldots,n\}\).
Equation~\eqref{eq:record-renewal-lower} implies
\[
 \liminf_{n\to\infty}
 \frac{E_{\widehat P_v}N_{\mathrm{rec}}(n)}{n}\ge\xi p_v>0.
\]
If \(E_{\nu_v}S=\infty\), however, the partial sums of iid copies
of \(S\), divided by their number, tend to infinity almost surely.
To see this, apply the strong law to \(\min(S,K)\) for each positive
integer \(K\), then let \(K\to\infty\).  Inverting these partial
sums shows \(N_{\mathrm{rec}}(n)/n\to0\) almost surely.  Since this
ratio lies in \([0,1]\), bounded convergence gives the same limit in
expectation, a contradiction.  Thus \(E_{\nu_v}S<\infty\), and
\(L\le S\) proves \eqref{eq:finite-mean-width}.
\end{proof}

The estimate concerns spatial width; it uses no moment assumption on
the time spent in a word.  We next show that coexistence of the two
escape signs also forces a finite mean for the full spatial radius.
